\section{Related work}
\label{sec:related}

\subsection{Log-based anomaly detection}

A mature line of work uses log templates and sequences to flag anomalous time windows. \textsc{DeepLog}~\cite{du2017deeplog} and \textsc{LogAnomaly}~\cite{meng2019loganomaly} train recurrent and Transformer-based models over per-line log templates extracted with Drain~\cite{he2017drain}, predicting whether the next log line is expected. \textsc{PLELog}~\cite{yang2021plelog}, \textsc{LogBERT}~\cite{guo2021logbert}, and \textsc{LogADempirical}~\cite{le2023logadempirical} extend this with BERT-style masked language modeling on log templates. These systems answer ``is something wrong'' --- the binary anomaly question. They do not answer ``what is wrong'' --- they do not retrieve a past similar incident. Our characterization confirms that telemetry-only models already saturate the binary anomaly task in our setting (PR-AUC 0.77 with a gradient-boosted classifier on 94 numeric features), so the marginal value of memory is in \emph{citation}, not detection.

\subsection{Multi-channel and trace-based diagnosis}

\textsc{Pinpoint}~\cite{chen2002pinpoint} and \textsc{TraceAnomaly}~\cite{liu2020traceanomaly} use distributed traces to localize root-cause services. \textsc{MicroRCA}~\cite{wu2020microrca} and \textsc{MEPFL}~\cite{zhou2019mepfl} jointly use traces and metrics for microservice fault localization. \textsc{Eadro}~\cite{lee2023eadro} integrates logs, traces, and metrics for end-to-end anomaly detection on microservice systems. These systems share our multi-channel approach but localize a \emph{service}, not a \emph{past Jira ticket}. The retrieval-versus-localization framing is different: localization tells you ``the bug is in cartservice''; retrieval tells you ``this incident matches INCIDENT-1234 --- cart Redis OOM, fixed by raising \texttt{maxmemory}.''

\subsection{Information retrieval over operational data}

\textsc{LogClass}~\cite{meng2020logclass} and \textsc{LogPunk}~\cite{liu2022logpunk} build classifiers over log signatures for category-level routing. \textsc{LogQA}~\cite{wang2023logqa} treats log analysis as question-answering over log text. The closest prior work to ours is \textsc{iLog}~\cite{jin2022ilog}, which retrieves past incident reports given a current log snippet using TF-IDF and ELMo embeddings. We extend that line of work in three ways: \emph{(a)} we use a cross-encoder reranker, which the prior work predates; \emph{(b)} we measure retrieval quality as a function of \emph{deployment-history depth} rather than treating the corpus as static; \emph{(c)} we honestly characterize the trade-off with anomaly detection (which prior work generally conflates with retrieval).

\subsection{Retrieval-augmented generation and dense retrieval}

\textsc{REALM}~\cite{guu2020realm}, \textsc{RAG}~\cite{lewis2020rag}, and \textsc{Atlas}~\cite{izacard2022atlas} framed retrieval as augmenting generative models with external knowledge. \textsc{ColBERT}~\cite{khattab2020colbert} introduced late-interaction retrieval bridging bi-encoder and cross-encoder approaches. \textsc{DPR}~\cite{karpukhin2020dpr} demonstrated dense retrieval beating BM25 on open-domain question answering. These methods inform our skill chain. The cross-encoder reranker we fine-tune is the same architectural class as MS-MARCO MiniLM~\cite{wang2020minilm}; we adapt it from generic-web ranking to the operational-incident domain. Our finding that the fine-tuned reranker shifts probability mass into Hit@5 at the cost of top-1 is consistent with the cross-encoder literature: joint scoring improves recall but is harder to calibrate at the top rank.

\subsection{Jira analytics}

\textsc{TAWOS}~\cite{tawosi2023tawos} released 458{,}000 real Jira issues with rich metadata. We use \textsc{TAWOS} as the empirical anchor for our V2 humanized corpus's length, comment-count, and code-block distributions; we do not train on \textsc{TAWOS}. \citet{diamantopoulos2018jira} extracted bug-report duplicates with topic modeling. Bichain et al.~\cite{bichain2023effortest} studied effort estimation from Jira descriptions. Our work is closer to the \emph{real-time retrieval} setting than to the \emph{post-hoc analysis} setting these papers occupy.

\subsection{What makes this paper distinct}

To our knowledge, no prior work has: \emph{(i)} empirically characterized retrieval quality as a function of \emph{deployment history depth} on a multi-channel SRE dataset; \emph{(ii)} identified the standard $\text{Recall@}K$ definition as misleading in the deep-history regime and recommended Hit@$K$ instead; \emph{(iii)} fine-tuned a cross-encoder on (telemetry-window, Jira-ticket) pairs and quantified the Hit@5-versus-Hit@1 trade-off; or \emph{(iv)} characterized cold-start novelty failures alongside the depth-scaling success, motivating future work on calibrated novelty heads.
